\section{Knowledge Does the Work}
\label{sec:results}
\label{sec:knowledge}

\subsection{Knowledge against scaffolding, on six of the seven models}
\label{sec:decomposition}

\begin{table}[t]
\centering
\footnotesize
\setlength{\tabcolsep}{2.6pt}
\begin{tabular}{@{}lrrrrr@{}}
\toprule
\tabhead{Model} & $k$ & \armS{} & \armH{} & \armM{} & \armC{} \\
\midrule
\mSol{} & 3 & 38.0 & 29.4 & 59.5 & \textbf{79.4} \\
\mQwen{}, official & 3 & 13.9 & 8.3 & 20.2 & \textbf{42.9} \\
\mQwen{}, end-to-end & 3 & 18.4 & 18.6 & 39.2 & \textbf{62.3} \\
\midrule
\mGLM{} & 1 & 13.9 & 19.3 & 22.9 & \textbf{55.1} \\
\mDS{} & 1 & 22.6 & 22.6 & 42.9 & \textbf{56.6} \\
\mSON{} & 1 & 22.9 & 38.0 & 23.8 & \textbf{68.4} \\
\mGPT{} & 1 & 37.0 & 51.2 & 50.3 & \textbf{77.4} \\
\bottomrule
\end{tabular}
\caption{DABStep hard-task accuracy (\%), $n{=}332$ (\mDS{}: 226 of its
complete-case set). The $k{=}3$ panels are the primary decomposition; the
four $k{=}1$ runs give breadth. \armS{} to \armH{} adds the governed tools,
the retrieval instruction and the rules with every description and SQL
definition emptied; \armH{} to \armC{} restores them. \mQwen{}'s
end-to-end row also credits an answer given as its final paragraph
(Section~\ref{sec:grading}).}
\label{tab:decomposition}
\end{table}

The four-arm ablation separates the contract's knowledge from the
scaffolding that delivers it. \armH{} keeps the contract's tools, its
instruction to look things up before writing SQL, its table allow-list
and its operation rules, and empties every description and SQL definition
(Section~\ref{sec:hollow}). Table~\ref{tab:decomposition} gives the
hard-task accuracy of all four arms on six models: two panels repeated
three times, and four single runs. (\mLuna{}, the
seventh model, did not run the hollow arm.)

\textbf{The knowledge step is large and certain on every model that ran
it.} Restoring
the knowledge (\armH{} to \armC{}) adds 50.0 hard-task points on \mSol{} (205
tasks to 12 over the three repeats; 95\% interval over the 38 question
shapes $+32$ to $+66$, 23 shapes to 5) and 34.5 to 43.8 on \mQwen{}
depending on the scoring view (189 tasks to 10 and to 8; over shapes, from
$+24$ at the lowest). On the four single runs it adds 26.2 to 35.8 points.
It is the fee procedure that it adds: on the 294 hard tasks that concern
fees the step is $+57.6$ points on \mSol{}; on the other 38 it is $-8.8$
on \mSol{} and $-0.9$ on \mQwen{} (end-to-end), and on the easy tasks $+13$
and $+12$.

\textbf{The scaffolding step is small, and its sign depends on the
model.} In the single runs it is $+0.0$ points on \mDS{}, $+5.4$ on \mGLM{}
and $+14$ to $+15$ on \mSON{} and \mGPT{}. Repeated, it is negative on the
newer frontier model: on \mSol{} the empty tools cost 8.5 hard-task points
(23 tasks to 66; over shapes $-16$ to $-2$, 5 shapes to 15), and the
deficit is significant in each repeat separately. On \mQwen{} the step is
$-5.5$ hard-task points under the official verdicts (over shapes $-10.5$ to
$-1.4$, though the shapes split 6 to 8) and $+0.2$ under the end-to-end
view ($p{=}0.4$), which also credits an answer given as the
last paragraph after its working (Section~\ref{sec:grading}). In the single runs the transcripts show how
empty tools can hurt: the agent calls \texttt{lookup\_domain}, is told
nothing is available, and searches harder; on \mGLM{} and \mSON{} \armH{}
burns more reasoning than any other arm (\appref{app:efficiency}). What every run agrees on is the
ordering of the two steps. \textbf{The knowledge step dominates the
scaffolding step on all six models}, and the sceptical reading that the
contract helps by narrowing the tables, imposing a procedure or slowing
the agent down predicts the opposite.

The manual pasted into the prompt carries the same domain, and it helps
too: \armM{} beats \armS{} by 21.6 hard-task points on \mSol{} (112 tasks
to 27) and by 6.3 or 20.8 on \mQwen{}. The contract beats the manual by
19.9 hard-task points on \mSol{} (116 tasks to 9; over shapes $+7$ to
$+34$) and 22.7 to 23.2 on \mQwen{}.
That gap looks like delivery, and Section~\ref{sec:delivery} takes it
apart.

\subsection{On LiveSQLBench, knowledge helps where it is determined}
\label{sec:lsb-knowledge}

LiveSQLBench supplies each database's knowledge base (KB), so the knowledge
arms there are built from the benchmark's own documentation
(Section~\ref{sec:benchmarks}). On the 302 primary tasks the KB as a manual
beats the bare schema on both models, as registered in (1): 35.5\%
against 20.0\% on \mLuna{} (71 tasks to 8, $p{<}10^{-10}$) and 37.3\%
against 23.2\% on \mQwen{} (68 to 13). The ceiling is low, and a blind test
says why.

Two model raters (Section~\ref{sec:threats}), each alone and blind to every
run and contract, read each task's question, gold SQL and result and the
KB entries the gold uses, and judged whether the question and the KB settle
every choice that changes the result. Each assigns one of four labels (the
KB is clear, ambiguous, contradicted by the gold, or silent on what the
gold needs); they agree on 83\% of tasks, $\kappa{=}0.79$ on settled or
not. Where both say settled, the task is \emph{determined}: 138 of 302.
Where both say not, it is \emph{undetermined}: 132; the raters split on
the other 32. The knowledge arms solve 60\% of determined tasks and 14\% of
undetermined ones, on each model alike. Because the raters saw the gold,
``determined'' partly means ``a reader of the KB would reach this gold'',
and determined tasks are also easier without knowledge (\armS{} 39\%
against 6\%). The knowledge's own gain over \armS{} is $+21$ points where
the KB settles the gold and $+8$ where it does not (Mann--Whitney on
per-task gains, $p{=}0.002$, \emph{expl.}), but from such different
baselines the comparison depends on the scale. In points the difference is
$+13$ (95\% interval over databases $+5$ to $+22$); in log-odds the two
gains are equal ($-0.04$, interval $-0.70$ to $+0.51$), and the 32 tasks the
raters split on gain the most. What the ratings show robustly is where the
accuracy is, not that knowledge works differently there. The raters call the KB
ambiguous on 109 tasks and contradicted by the gold on 43 to 45; in a
closer reading of 55 tasks the open choices were most often how a
\texttt{NULL} or an unlisted category is scored, which fields a formula
averages, and which of two columns holds a quantity.
Section~\ref{sec:authoring} returns to what an author can and cannot add
there.

\subsection{What a pre-computed layer would add}
\label{sec:ceiling}

Compiling the DABStep contract's own SQL expressions into two views gives
an oracle for the layer a context-layer product would ship; no agent is
given the views. Instantiated by us with each task's parameters, they
reproduce the gold answer on every one of the 176 hard tasks whose
question shape they cover, selected by a pattern over the question text.
On those tasks whatever an agent misses is a failure to derive from
knowledge it provably had. In the four single runs that \emph{derivation
gap} is 39.2 points on \mGLM{}, 37.3 on \mDS{}, 26.1 on \mSON{} and 5.1 on
\mGPT{}, falling with capability (single runs, so a trend, not a
per-model estimate), while
the shortfall on the 148 tasks that need reasoning beyond the views moves
under no ordering (\appref{app:buckets}). Pre-computing a metric pays
where agents fail to derive it, and on the strongest of the four that is
5 points of the tasks the views cover. The contract's gain also stays
inside its domain: on the hard tasks that do not concern fees it holds no
detectable advantage in any of the four runs (\appref{app:domain}).
